\subsection*{58. Make the Human-Feedback Budget Visible}

Suppose a pilot has \textbf{12 learners}. Each receives 4 substantive artifact reviews averaging \textbf{25 minutes}, plus 1 \textbf{15-minute} technical defense. That is:
\begin{equation}
\begin{aligned}
12 \times (4 \times 25 + 15)\text{ minutes} &= 1{,}380\text{ minutes} \\
&= 23\text{ reviewer-hours}.
\end{aligned}
\end{equation}

Add preparation, live instruction, review calibration, administration, marketing, learner support, retries, and contingencies. If the founder has 15 to 25 total sustainable hours a week alongside a full-time job, even a seemingly small 4-week program can consume a meaningful share of available capacity. The institution must not sell ``unlimited mentorship'' while privately hoping questions will be rare.

2 ledgers are necessary: actual cash expenditure during the experiment and \textbf{normalized delivery economics} including the market cost of the founder's skilled labor. A course can be affordable to pilot because the founder personally supplies labor without taking a salary, yet fail to become a repeatable business when equivalent work must be paid for.

\subsection*{59. A Learner Journey Is Also an Error-Recovery Journey}

Before purchase, an applicant should see prerequisites, sample instruction, the expected time commitment, actual teacher identity, assessment conditions, and refund or withdrawal terms. On admission, the learner receives a diagnostic and a realistic study schedule. Each week contains preparation, live practice, feedback, and revision. If the learner misses work, a defined recovery path identifies whether the obstacle is time, access, confidence, or a prerequisite gap.

At the end, the institution distinguishes participation from independent achievement. Some learners may need a further attempt; some may not meet the standard. A trustworthy program can celebrate effort without mislabeling mastery. Learners should receive a useful capability report, including limitations, rather than a grand certificate whose meaning cannot be defended.

\subsection*{60. The College-Partnership Architecture Is Genuinely Different}

A curriculum-integrated course in a college may involve a larger and more heterogeneous cohort, fixed academic timetables, institutional faculty, examination policies, student data systems, and formal approval processes. The most valuable Bayes contribution may be curriculum architecture, instructor preparation, authentic labs, assessment rubrics, and quality assurance, not direct founder teaching for every student.

A viable pilot should specify which faculty teach, which parts are jointly designed, who grades, how markers calibrate, how student access is managed, and what institutional authorization exists. It must resist the temptation to promise individualized boutique mentoring for hundreds of students without the supporting workforce. Evaluate the \emph{faculty-enabled system}, not merely the quality of a visiting masterclass.

\subsection*{61. The Enterprise Architecture Begins with Work, Not Course Topics}

An enterprise buyer may care less about understanding the history of AI architectures than about improving a controlled workflow. The learning target should be derived from an actual task: for instance, building a reliable internal evaluation procedure or safely using AI to reduce repetitive analysis without exposing confidential data. The manager must provide protected time, data access where appropriate, and an opportunity to apply the skills.

Delivery might combine a short baseline task, domain-specific examples, supervised application to approved work, and a measurable post-training implementation. A polished workshop is not enough if the workplace refuses to adopt the new process. The instructional design therefore includes \textbf{transfer conditions outside the classroom}: management support, workflow integration, incentives, and permission to change how work is done.

\subsection*{62. What Bayes Can and Cannot Claim After the 1st Experiment}

If learners complete unfamiliar tasks to defined standards, Bayes may responsibly describe those observed results with sample sizes, conditions, and limitations. If some learners improve between entry and exit tasks, it may describe that pre-post change. It cannot automatically attribute every gain to instruction or claim salary and employment outcomes. A small success does not prove all adult engineers will benefit or that a 6-month flagship is commercially viable.

The more valuable result of the 1st pilot may be the \textbf{institutional learning}: a validated teaching sequence for 1 skill, a realistic feedback-minute budget, a working rubric, a map of predictable learner errors, and evidence of which kind of buyer values the experience. These are assets an early institute can compound before it possesses a campus or a famous brand.
